\begin{ExerciseBlock}
\ExerciseHeading{03}{Trägerrakete stößt Raumkapsel ab}{two_particle_systems_space_capsule_separation}
Die letzte Stufe einer Trägerrakete hat die Gesamtmasse $M = m_1 + m_2 = 1000\,$kg und fliegt mit konstanter Geschwindigkeit $v_0 = 12000\,$m/s geradlinig durch den Weltraum. Sie stößt eine Raumkapsel der Masse $m_2 = 200\,$kg mit der Relativgeschwindigkeit $v_\mathrm{rel}=\qty{1.0}{\metre\per\second}$ in Flugrichtung ab.

\begin{center}
	\CachedTikzFigure{content/resources/chapters/two_particle_systems/tikz/space_capsule_separation.tex}
\end{center}

\begin{enumerate}
	\item Mit welchen Geschwindigkeiten fliegen die Raketenstufe und die Kapsel anschließend weiter?

	\item Wie viel Energie muss die Trägerrakete beim Abstoßen der Kapsel aufbringen und in kinetische Energie überführen? Handelt es sich hierbei um Wärmeenergie? Was unterscheidet diesen Prozess von einem inelastischen Stoß?
\end{enumerate}
\end{ExerciseBlock}
\UseExerciseSolution{two_particle_systems_space_capsule_separation}
